\documentclass[letterpaper, 11pt]{article}
\usepackage{latexsym}
\usepackage[empty]{fullpage}
\usepackage{titlesec}
\usepackage{marvosym}
\usepackage[usenames, dvipsnames]{color}
\usepackage{verbatim}
\usepackage{enumitem}
\usepackage[hidelinks]{hyperref}
\usepackage{fancyhdr}
\usepackage[english]{babel}
\usepackage{tabularx}
\usepackage{fontawesome5}
\usepackage{multicol}
\setlength{\multicolsep}{-3.0pt}
\setlength{\columnsep}{-1pt}
\input{glyphtounicode}
\usepackage{vwcol}
\usepackage{comment}
\usepackage{hyperref}
\hypersetup{
	colorlinks=true,
	linkcolor=black,
	filecolor=black,
	citecolor=black,
	urlcolor=Blue,
}

%----------FONT OPTIONS----------
% sans-serif
% \usepackage[sfdefault]{FiraSans}
% \usepackage[sfdefault]{roboto}
% \usepackage[sfdefault]{noto-sans}
% \usepackage[default]{sourcesanspro}

% serif
% \usepackage{CormorantGaramond}
% \usepackage{charter}

\pagestyle{fancy}
\fancyhf{} % clear all header and footer fields
\fancyfoot{}
\renewcommand{\headrulewidth}{0pt}
\renewcommand{\footrulewidth}{0pt}

% Adjust margins
\addtolength{\oddsidemargin}{-0.6in}
\addtolength{\evensidemargin}{-0.5in}
\addtolength{\textwidth}{1.19in}
\addtolength{\topmargin}{-.5in}
\addtolength{\textheight}{1.1in}

\urlstyle{same}

\raggedbottom
\raggedright
\setlength{\tabcolsep}{0in}

% Sections formatting
\titleformat{\section}{ \vspace{-4pt}\scshape\raggedright\large\bfseries }{}{0em}{}[
\color{black}
\titlerule
\vspace{-5pt}
]

% Ensure that generate pdf is machine readable/ATS parsable
\pdfgentounicode=1

%-------------------------
% Custom commands
\newcommand{\resumeItem}[1]{ \item\small{ {#1 \vspace{-2pt}} } }
%\newcommand{\classesList}[4]{ \item\small{ {#1 #2 #3 #4 \vspace{-2pt}} } }

\newcommand{\resumeSubheading}[4]{
\vspace{-1pt}
\item
\begin{tabular*}{1.0\textwidth}[t]{l@{\extracolsep{\fill}}r}
	\textbf{#1}       & {\small #2} \\
	{\small#3} & \textit{\small #4} \\
\end{tabular*}
\vspace{-10pt}
}

\newcommand{\resumeProjectHeading}[2]{
\vspace{-1pt}
\item
\begin{tabular*}{1.0\textwidth}[t]{l@{\extracolsep{\fill}}r}
	\textbf{#1}  & {\small #2} \\
\end{tabular*}
\vspace{-20pt}
}

% \newcommand{\resumeSubSubheading}[2]{ \item
% \begin{tabular*}{0.97\textwidth}{l@{\extracolsep{\fill}}r}
% 	\textit{\small#1} & \textit{\small #2} \\
% \end{tabular*}
% \vspace{-7pt}
% }

\newcommand{\resumeSubItem}[1]{\resumeItem{#1}
\vspace{-4pt}}

\renewcommand{\labelitemi}{$\vcenter{\hbox{\tiny$\bullet$}}$}
\renewcommand{\labelitemii}{$\vcenter{\hbox{\tiny$\bullet$}}$}

\newcommand{\resumeSubHeadingListStart}{\begin{itemize}[leftmargin=0.0in, label={}]}
\newcommand{\resumeSubHeadingListEnd}{\end{itemize}}
\newcommand{\resumeItemListStart}{\begin{itemize}[leftmargin=0.25in]}
\newcommand{\resumeItemListEnd}{\end{itemize}
\vspace{-5pt}}

%-------------------------------------------
%%%%%%  RESUME STARTS HERE  %%%%%%%%%%%%%%%%%%%%%%%%%%%%

\begin{document}
	
%----------HEADING----------
\begin{center}
    {\Huge \scshape Gawtam Chithra Ramesh} \\
    \vspace{3pt}
    \small
    { \href{tel:+46764518589}{{\faMobile*} \textcolor{black}{(+46) 76541-8589}} ~
    \href{mailto:gawtamcr3@gmail.com}{{\faPaperPlane} \textcolor{black}{gawtamcr3@gmail.com}} ~
    \href{https://linkedin.com/in/gawtamcr/}{{\faLinkedinIn} \textcolor{black}{Gawtam C R}} ~
    \href{https://github.com/gawtamcr}{{\faGithub} \textcolor{black}{gawtamcr}} ~
    \href{https://gawtamcr.github.io/}{{\faGlobe} \textcolor{black}{Portfolio}} ~
    \href{https://maps.google.com/?q=Stockholm,Sweden}{{\faMapMarker*} \textcolor{black}{Stockholm, Sweden}}
    }
    \vspace{-8pt}
\end{center}

%-----------SUMMARY-----------
\section{Summary}
\small{Machine-learning researcher (MSc Systems, Control \& Robotics, KTH) working on \textbf{reinforcement learning} and \textbf{end-to-end learning} for embodied agents, with reinforcement learning the through-line of my work since 2020. My experience spans \textbf{multi-agent PPO} policies for games and a \textbf{generative flow-matching} trajectory planner in my Master's thesis at ABB Robotics. Comfortable across \textbf{generative modeling}, deep and multi-agent reinforcement learning, and \textbf{motion planning and control}, with an independent, rapid-prototyping research style and three first-author publications.}
\vspace{-4pt}

%-----------EDUCATION-----------
\section{Education}
    \resumeSubHeadingListStart 
    \resumeSubheading {KTH Royal Institute of Technology}{Aug 2024 -- Jul 2026}{Master of Science in \href{https://www.kth.se/en/studies/master/systems-control-robotics/msc-systems-control-and-robotics-1.8733}{Systems, Controls, and Robotics} (GPA 4.62/5.0)}{Stockholm, Sweden}
    
    \resumeSubheading {Indian Institute of Technology Madras}{Aug 2019 -- Jul 2024}{Dual Degree (B.Tech in \href{https://ed.iitm.ac.in/about.html}{Engineering Design} and M.Tech in \href{https://ed.iitm.ac.in/~robotics_lab/}{Robotics}) (GPA 7.9/10)}{Chennai, India}

    \resumeSubHeadingListEnd

%-----------PUBLICATIONS-----------
\section{Publications}
\begin{itemize}[leftmargin=0.15in]
    \setlength{\itemsep}{2pt}
    \item ``Robust Object-layer Construction of 3D Scene Graphs Using Instance Segmentation.'' \textit{International Conference on Image Processing and Vision Engineering (IMPROVE 2026)}. (Accepted, Long Oral, Best Paper Award)
    \item ``Synchronized RGB-D Inpainting for Privacy-Aware 3D Scene Graph Construction.'' \textit{ICPR 2026 Workshop on Perception, Interaction and Decision-Making for Human Cyber-Physical Systems (PIDM4HCPS)} . (Full paper, Accepted)
    \item ``Scene Understanding in Deformable Object Manipulation via Taxonomy-Guided Vision-Language Models.'' \textit{IROS 2025 Workshop on Robotic Manipulation of Deformable Objects (ROMADO)}. (Short paper, Accepted)
\end{itemize}

%-----------EXPERIENCE-----------
\section{Professional Experience}
\resumeSubHeadingListStart

    \resumeSubheading{Master's Thesis on Generative Planning for Long-Horizon Manipulation}{Jan 2026 -- Jun 2026 (Ongoing)}{\textbf{ABB Robotics}, Mentors:
    \href{https://www.linkedin.com/in/matthewwilliamlock/}{Matthew Lock} and \href{https://www.linkedin.com/in/jonathan-styrud-99385864/?locale=en_US}{Jonathan Styrud}}{Västerås, Sweden} 
    \resumeItemListStart 
        \resumeItem{Building a \textbf{generative flow-matching trajectory planner} that learns a motion prior from offline demonstrations and is steered at \textbf{inference time} toward new objectives, with no retraining or per-task expert data.}
        \resumeItem{Conditioning one model on objectives at test time so a single planner transfers across task variants, rather than training a separate \textbf{imitation-learning} policy per task.}
        \resumeItem{Using \textbf{differentiable specification gradients} (Signal Temporal Logic) to steer generation toward compositional, temporally-extended goals.}
        \resumeItem{Validating on \textbf{long-horizon planning benchmarks} against state-of-the-art baselines, and transferring the approach to \textbf{6-DoF manipulation} with ABB Robotics R\&D.}
    \resumeItemListEnd
    
    \resumeSubheading{Device Research Intern on 3D Scene Graph}{Jul 2025 -- Dec 2025}{\textbf{Ericsson} Research, Mentors:
    \href{https://se.linkedin.com/in/p\%C3\%BCren-g\%C3\%BCler-76457317}{{Puren Guler}} and 
    \href{https://se.linkedin.com/in/hcaltenco}{Hector Caltenco}}{Lund, Sweden} 
    \resumeItemListStart 
        \resumeItem{Developed a robust 3D scene-graph pipeline by integrating \textbf{instance segmentation and tracking} into a SOTA framework (Hydra), reducing sensitivity to parameter tuning across diverse environments.}
        \resumeItem{Designed a synchronized \textbf{RGB-D inpainting} module for privacy-aware 3D scene-graph construction, maintaining consistency across RGB and depth channels.}
        \resumeItem{Analyzed 3DSG frameworks to identify research gaps in object merging, parameter sensitivity, and privacy protection.}
        \resumeItem{Benchmarked SOTA inpainting algorithms to balance reconstruction fidelity with real-time performance.}
    \resumeItemListEnd

    \resumeSubheading{Graduate Robotics Intern on Modular Collaborative Robots}{Dec 2022 -- Jul 2023}{\textbf{Systemantics} India Pvt. Ltd., Mentor: \href{https://www.linkedin.com/in/jagannathraju/}{{Jagannath Raju}}}{Bangalore, India}
    \resumeItemListStart 
        % \resumeItem{Implemented backward-chained Behavior Trees for long-horizon 6DOF manipulator planning using ROS2 and Moveit2.}
        % \resumeItem{Implemented a control synthesis pipeline for a 6DOF manipulator using Behavior Trees to enable robust, long-horizon trajectory planning.}
        % \resumeItem{Designed and implemented a long-horizon trajectory planning system for a 6DOF manipulator using Behavior Trees,}
        \resumeItem{Developed backward-chained \textbf{Behavior Trees} for a 6-DoF manipulator to enable robust, long-horizon trajectory planning.}
        \resumeItem{Designed and implemented a 6-DoF \textbf{trajectory generation and control} pipeline in C++, achieving \textbf{real-time} joint actuation via low-latency communication with motor controllers using D-Bus and SocketCAN.}
        % \resumeItem{Developed a S-Curve trajectory generator for the arm, streaming real-time joint commands directly to motor controllers via DBUS and SocketCAN.}
        % \resumeItem{Developing an disturbance observer and admittance control model for all the 6 joints of the robotic arm}
    \resumeItemListEnd
    \resumeSubheading{Research Intern, Multi-Object Tracking}{Jun 2021 -- Jul 2021}{Blurgs Research Labs Pvt. Ltd.}{Remote, India}
    \resumeItemListStart
        \resumeItem{Analyzed 5 state-of-the-art multi-object tracking networks for the startup's first product for the Indian Navy.}
        \resumeItem{Trained and tested multi-object tracking networks (\textbf{Siam-MOT}, \textbf{FairMOT}) for surveillance with drones and CCTV.}
        \resumeItem{Used \textbf{FairMOT} to track and trail multiple moving objects and re-identify them in real time.}
    \resumeItemListEnd

    \newpage
    % \resumeProjectHeading{Graduate Teaching Assistant $|$ \textnormal{\small Field and Service Robotics (ID4060 - IITM)} } {Jan 2025 -- Feb 2025}
    \resumeSubheading {Graduate Teaching Assistant}{Aug 2023 -- Nov 2023}{Field and Service Robotics (ID4060), Mentor: \href{https://bijosebastian.wordpress.com/}{Bijo Sebastian}}{IIT Madras, India}
    \resumeItemListStart 
        \resumeItem{Involved in design and evaluation of assignments based on ROS and CoppeliaSim for 40 students.}
        \resumeItem{Held meetings with small groups of students with diverse backgrounds to review their assignments and projects.}
    \resumeItemListEnd

    \resumeSubheading{Strategist, Mentor \& Co-organizer}{Apr 2020 -- May 2021}{Computer Vision and Intelligence Group (CVIG), IIT Madras}{Chennai, India}
    \resumeItemListStart
        \resumeItem{Co-organized a deep-learning summer school for 200+ students spanning \textbf{deep learning}, \textbf{reinforcement learning}, and \textbf{computer vision}.}
        \resumeItem{Co-mentored student projects on \textbf{reinforcement learning} for games, an early thread of the RL focus that runs through my work today.}
    \resumeItemListEnd

\resumeSubHeadingListEnd
%-----------EXPERIENCE-----------
\section{Research Experience}
\resumeSubHeadingListStart
 
    \resumeSubheading{Graduate Robotics Research on Deformable Object Manipulation}{Jan 2025 -- Nov 2025}{\href{https://www.kth.se/is/rpl}{Robotics Perception and Learning}, Advisor:
    % \href{https://www.kth.se/is/rpl}{David Blanco-Mulero} 
    %\href{https://yifeidong0.github.io/}{\emph{Yifei Dong}} and 
    \href{https://www.kth.se/is/rpl}{Florian T. Pokorny}}{KTH, Sweden} 
    \resumeItemListStart 
        % \resumeItem{Deformable object manipulation planning through taxonomy-based abstractions and foundation models}
        % \resumeItem{Developed a framework integrating pretrained-VLMs with a structured taxonomy to enhance DOM.}
        % \resumeItem{Conducted in-depth analysis of VLM performance, identifying challenges in DOM scene understanding and proposing solutions like refined taxonomy and multi-modal inputs; }
        % \resumeItem{Submitted and presented the result at ROMADO, a \textbf{IROS2025 Workshop}.}
        % \resumeItem{Designed and implemented a simulation pipeline for task execution using PyBullet, with plans for real-world validation 
        % %to assess sim-to-real transferability.
        % }

        % \resumeItem{Integrated taxonomy-guided vision-language models (Gemini, Qwen) to interpret DOM scenes and map 
        % structured outputs to 
        % motion primitives}
        % \resumeItem{Built prompt–taxonomy framework to generate structured codes and textual justifications for robot-parsable perception}
        \resumeItem{Integrated taxonomy-guided \textbf{vision-language models} (Gemini, GPT-4o, Qwen) to interpret deformable-object scenes and map their structured outputs to motion primitives.}
        \resumeItem{Built a prompt and taxonomy framework generating structured, robot-parsable specifications from natural-language instructions for downstream manipulation.}
        \resumeItem{Performed quantitative failure-mode analysis to guide iterative refinement of the taxonomy and prompts, integrating multimodal and temporal cues.}
    \resumeItemListEnd

    \resumeSubheading{Dual Degree Project on Structure from Motion}{Jan 2024 -- May 2024}{\href{https://niravatgit.github.io/INSPIRELab/}{INSPIRE Lab}, Advisor: \href{https://niravatgit.github.io/patelnirav.com/index.html}{{Nirav Patel}}}{IIT Madras, India}
    \resumeItemListStart 
        \resumeItem{Developed a \textbf{3D perception} system using COLMAP (\textbf{Structure from Motion}) from monocular RGB images to build cost-efficient human torso reconstruction for more accessible ultrasound.}
        \resumeItem{Developed an autonomous ultrasound scanning system using ROS, integrating a 5-DoF robotic arm with a single RGB camera to achieve cost-effective 3D perception for medical imaging.}
        \resumeItem{Implemented motion planning and control using ROS and MoveIt, and simulated the 5-DoF manipulation in Gazebo.}
        \resumeItem{Integrated perception algorithms within ROS, processing 3D point clouds from RGB images to enable accurate trajectory planning for the ultrasound probe.}
    \resumeItemListEnd

    \resumeSubheading{Offline And Online Motion Planning \& Control of Unmanned Aerial Vehicles}{Sep 2021 -- Apr 2022}{\href{https://yrf.iitm.ac.in/}{Young Research Fellowship}, Advisor: \href{http://www.ae.iitm.ac.in/~satadal/index.html}{\emph{Prof. Satadal Ghosh}}}{IIT Madras, India}
    \resumeItemListStart
        \resumeItem{Developed a \textbf{target-driven motion planning} and \textbf{obstacle avoidance} algorithm for UAVs.}
        \resumeItem{Used Proportional Navigation and Acceleration-Velocity Obstacles in a 2D field to achieve a 20\% faster and more compute-efficient algorithm than classical aerospace methods.}
    \resumeItemListEnd

\resumeSubHeadingListEnd



%-----------PROJECTS-----------
\section{Selected Projects}
\resumeSubHeadingListStart

   \resumeProjectHeading{Deep Learning, Advanced: Generative \& Self-Supervised Models $|$ \textnormal{\small DD2610, KTH}}{Nov 2025 -- Jan 2026}
    \resumeItemListStart
        \resumeItem{Reimplemented \textbf{MeanFlow}, a one-step \textbf{generative flow-matching} model, from scratch in \textbf{PyTorch}; ablated time samplers, classifier-free guidance ($\omega, \kappa, p_c$), and the $r\neq t$ ratio, reaching best FID $\approx$160 on CIFAR-10 (epoch 600) and 75\% optimal $r{=}t$ mixing on MNIST -- directly informing my thesis on generative trajectory planning.}
        \resumeItem{Built \textbf{Masked Autoencoders} with a \textbf{Vision Transformer} encoder from scratch (75\% patch masking, MSE on masked patches), reducing reconstruction loss from 0.385 to 0.127 over 5 epochs on CIFAR-10.}
        \resumeItem{Implemented \textbf{contrastive} (\textbf{SimCLR}, NT-Xent $\tau{=}0.07$) and \textbf{semi-supervised} (\textbf{FixMatch}, confidence threshold 0.95) learners in \textbf{JAX} with a ResNet-18 backbone, reaching $\sim$24\% linear-probe accuracy from only 1\% of labels.}
    \resumeItemListEnd

   \resumeProjectHeading{Robot Learning and Embodied AI $|$ \textnormal{\small DD2600, KTH}}{Sep 2025 -- Oct 2025}
    \resumeItemListStart 
        \resumeItem{Implemented a 3D semantic object mapping pipeline using RGB-D data, integrating \textbf{open-vocabulary object detection}, segmentation, and \textbf{CLIP} for semantic grounding.}
        \resumeItem{Deployed a pre-trained \textbf{Vision-Language-Action (VLA)} model for robot manipulation, executing \textbf{policy inference} from visual inputs and natural-language commands.}
        \resumeItem{Characterised generalisation and robustness failures of the VLA on unseen tasks, documenting the conditions under which it broke down.}
    \resumeItemListEnd
    \newpage
    \resumeProjectHeading{Multi Agent Path Finding in Unity3D $|$ \textnormal{\small DD2438, KTH}}{Feb 2025 -- May 2025}
    \resumeItemListStart 
        \resumeItem{Developed a \textbf{Hybrid A*} path planner and waypoint tracking system in Unity3D for kinematically constrained single-agent (car/drone) navigation, achieving competitive performance.}
        \resumeItem{Engineered \textbf{Conflict-Based Search (CBS)} to solve multi-agent path finding for over 20 holonomic (drones) and non-holonomic (cars) agents, minimizing makespan in shared environments. }
        \resumeItem{Implemented a greedy dynamic goal assignment strategy for multi-agent teams ("bakery problem") integrated with CBS, optimizing task allocation and minimizing overall travel time. }
        % \resumeItem{battle snake }
    \resumeItemListEnd

    \resumeProjectHeading{Reinforcement Learning for Pacman Capture The Flag $|$ \textnormal{\small DD2438, KTH}}{Mar 2025 -- May 2025}
    \resumeItemListStart 
        \resumeItem{Designed a hybrid AI for Pacman CTF (Unity3D), merging a custom \textbf{Behavior Tree} for high level strategy with \textbf{RL} models for specialized actions (attack, defend, escape).}
        \resumeItem{Trained \textbf{RL (PPO)} policies via \textbf{Unity ML-Agents} with role-specific designs (observations, rewards) and advanced methods (phased training, \textbf{self-play, curriculum learning}) for adaptive agent behavior.}
        \resumeItem{Evaluated strategic agent decision-making and adaptive behavior switching in a competitive, partially observable, dynamic \textbf{multi-agent} Pacman environment.}
    \resumeItemListEnd

    \resumeProjectHeading{Image analysis and Computer vision $|$ \textnormal{\small DD2423, KTH}}{Oct 2024 -- Nov 2024}
    \resumeItemListStart 
        \resumeItem{Applied \textbf{Fourier Transforms} for image filtering and frequency analysis; designed Gaussian smoothing filters via FFT to reduce noise and study effects on varying image resolutions.}
        \resumeItem{Implemented a multi-scale differential geometry edge detector (computing $L_{vv}, L_{vvv}$ for zero\-crossings)
        and used Hough Transform for line detection from edges.}
        \resumeItem{Implemented image matching (\textbf{SIFT}, \textbf{RANSAC}) to estimate homographies (planar scenes) and fundamental matrices, enabling 3D scene reconstruction via triangulation}
    \resumeItemListEnd

    \resumeProjectHeading{Introduction to robotics $|$ \textnormal{\small DD2410, KTH}}{Sep 2024 -- Oct 2024}
    \resumeItemListStart 
        \resumeItem{Implemented an iterative \textbf{inverse kinematics} solver from scratch for a 7-DoF KUKA robot arm in ROS.}
        \resumeItem{Applied \textbf{RRT*} to find collision-free paths for a \textbf{Dubins car}, showcasing path planning for non-holonomic vehicles.}
        \resumeItem{Engineered 2D \textbf{occupancy grid mapping} solutions in ROS, processing laser scans to build maps, identify free space, and generate C-spaces for navigation.}
        \resumeItem{Integrated a \textbf{Behavior Tree} mission planner for a TIAGo mobile manipulator (using ROS and Gazebo), enabling autonomous navigation, vision-based grasping, and robust task execution (e.g., kidnapped robot recovery).}
    \resumeItemListEnd
    
    % \resumeSubheading{Control of Automotive Systems}{Jun 2022 - Nov 2022}{\emph{Mentored by \href{https://home.iitm.ac.in/srikanthan/}{ Prof.Srikanthan Sridharan}}}{IIT Madras, India}
    \resumeProjectHeading{Control of Automotive Systems $|$ \textnormal{\small ED4060, IIT Madras}}{Jun 2022 -- Nov 2022}
    \resumeItemListStart
        \resumeItem{Designed a \textbf{heading-angle controller} for an \textbf{autonomous ground vehicle} from its \textbf{vehicle dynamics}, reducing error between the desired and actual heading angle within given performance specifications.}
        \resumeItem{Improved performance of a pneumatic brake system by designing \textbf{P and PI controllers} to meet given specifications.}
        \resumeItem{Designed a nonlinear \textbf{Sliding Mode Controller} to attenuate the effect of external forces on a tractor's hydraulic hitch, reducing excessive vibrations and preventing front-wheel lift-off.}
    \resumeItemListEnd

    \resumeProjectHeading{Multi-Agent Reinforcement Learning Shooter $|$ \textnormal{\small CFI, IIT Madras}}{May 2020 -- Aug 2020}
    \resumeItemListStart
        \resumeItem{Trained cooperative \textbf{multi-agent PPO} policies in Unity3D for a team-based shooter, learning to outplay human opponents.}
        \resumeItem{Integrated the trained agents into a VR build; one of my first \textbf{reinforcement-learning} systems and the start of a long-running interest in learned multi-agent behaviour.}
    \resumeItemListEnd

\resumeSubHeadingListEnd


%-----------SKILLS-----------
\section{Technical Skills}

\begin{multicols}{2}
    \begin{itemize}[itemsep=- 1pt, parsep=3pt]
    \small{ 
              \item \textbf{Languages} : Python, C, C++, C\texttt{\#}\\ 
              \item \textbf{Robotics} : ROS2, Unity3D, MATLAB, MuJoCo\\ 
              \item \textbf{DL} : PyTorch, JAX, OpenCV, wandb \\
              \item \textbf{Tools} : VS Code, Docker, Git, \LaTeX, CUDA, Slurm
        }
    \end{itemize}
\end{multicols}

%------RELEVANT COURSEWORK-------
\section{Relevant Coursework}

\textbf{Machine Learning}: Robot Learning \& Embodied AI, Deep Learning Advanced, Artificial Intelligence \& Multi-Agent Systems, Fundamentals of Deep Learning \\
\textbf{Robotics \& Vision}: Introduction to Robotics, Image Analysis \& Computer Vision, Applied Estimation, Field and Service Robotics \\
\textbf{Controls}: Modelling of Dynamical Systems, Control Theory Advanced \& Practice, Control of Automotive Systems \\

%-----------ACHIEVEMENTS---------------
\section{Achievements}
\begin{itemize}[leftmargin=0.15in]
    \setlength{\itemsep}{-0.3em}
    \item Awarded the IIT Madras Young Research Fellowship (top 30 of 800+) for outstanding research potential.
    \item JEE Advanced 97.64 percentile and JEE Mains 99.10 percentile (1.5M candidates); UCEED All-India Rank 52 of 50,000+.
\end{itemize}

\end{document}